%% core/c3-langlands.tex —— C3 轨道：Hecke 算子与特征层、单极子模空间与 Hecke 修正空间、
%% D-模与平坦联络、几何 Langlands 对应、翻译器 B_cc
%% 七段式：\oneline / definition / theorem+proof / \physside / example / \statusline / \dictionaryrow

\begin{centry}{Hecke 算子与 Hecke 特征层}{Hecke Operators and Hecke Eigensheaves}{hecke}
\oneline{在曲线的每一点按「一次丛的修正」定义的对应与算子；它是几何 Langlands 的算子侧，而它的本征对象——Hecke 特征层——正是与 ${}^{L}G$-局部系统配对的那一半词典。}

\begin{definition}[Hecke 对应与 Hecke 函子]
设 $C$ 为 $\mathbb{C}$ 上光滑射影连通曲线，$G$ 为复约化群，${}^{L}G$ 为其 Langlands 对偶群，$z\in C$。\emph{$z$ 处的 Hecke 对应}是
\[
\mathrm{Bun}_G\ \xleftarrow{\ p_z\ }\ \mathcal{H}_{z}\ \xrightarrow{\ q_z\ }\ \mathrm{Bun}_G ,
\]
其中 $\mathcal{H}_{z}$ 以三元组 $(E,E',\beta)$ 为点：$E,E'$ 为 $C$ 上 $G$-丛，$\beta:E|_{C\setminus\{z\}}\xrightarrow{\ \sim\ }E'|_{C\setminus\{z\}}$ 为 $z$ 外的同构；$p_z$ 记 $E$（忘记修正），$q_z$ 记 $E'$（修正后的丛）。纤维 $p_z^{-1}(E)$ 含于仿射 Grassmann 流形 $\mathrm{Gr}_{G,z}=G(\mathbb{C}(\!(t)\!))/G(\mathbb{C}[[t]])$；$G=GL_n$ 且只许一阶极点时它就是 $\mathbb{P}^{n-1}$。

设 $V$ 为 ${}^{L}G$ 的有限维表示。\emph{Hecke 函子}是
\[
H_{V,z}(\mathcal{F})\;=\;(q_z)_{*}\bigl(p_z^{*}\mathcal{F}\otimes\mathcal{L}_{V}\bigr),
\]
其中 $\mathcal{L}_{V}$ 是几何 Satake 对应给出的、沿修正方向的层；对 $G=GL_n$、$V=V_{\check\omega_1}$ 取平凡层，此时 $H_{V,z}(\mathcal{F})=(q_z)_{*}p_z^{*}\mathcal{F}$（Frenkel 2009 §2 式 (2.2)）。称 $D$-模 $\mathcal{F}$ 为\emph{Hecke 特征层}、以平坦 ${}^{L}G$-丛 $E$ 为特征值，若对每个 $z$ 与每个 $V$ 有
\[
H_{V,z}(\mathcal{F})\;\cong\;V_{E}|_{z}\otimes\mathcal{F},\qquad V_{E}=E\times_{{}^{L}G}V .
\]
\end{definition}

\begin{theorem}[Hecke 特征层的存在，$GL_n$]
设 $C$ 为 $\mathbb{C}$ 上光滑射影连通曲线，$n\ge1$，$G=GL_n$，$E$ 为 $C$ 上不可约的秩 $n$ 平坦丛。则存在非零 $D$-模 $\mathcal{F}_{E}$，使对一切 $z\in C$ 与 ${}^{L}G$ 的有限维表示 $V$ 满足 $H_{V,z}(\mathcal{F}_{E})\cong V_{E}|_{z}\otimes\mathcal{F}_{E}$。
\end{theorem}
\begin{proof}
先用几何 Satake 对应（Beilinson--Drinfeld、Mirkovi\'c--Vilonen）把 $E$ 与 ${}^{L}G$ 的表示丛转成仿射 Grassmann 流形上沿 Schubert 簇支撑的 perverse 层；再在 $\mathrm{Bun}_G$ 上对它作「逐点卷积」，得到 $D$-模 $\mathcal{F}_{E}$。卷积的结合律与 Hecke 对应的泛性质给出特征方程：沿每个 $z$ 作一次修正后推回，恰等于张量 $V_{E}|_{z}$。$GL_1$ 情形就是 Fourier 变换（Deligne；Laumon），$GL_2$ 由 Drinfeld 构造，一般 $n$ 由 Frenkel--Gaitsgory--Vilonen 完成。$E$ 不可约时取 $\mathcal{F}_{E}$ 的一个不可约成分即可。
\end{proof}

\physside{数学侧的 $H_{V,z}$ 是对 $\mathrm{Bun}_G$（或 $\mathcal{M}_H$）上的对象作一次 Hecke 修正；物理侧它是支撑在 $\mathbb{R}\times\{z\}$ 上的静态 't Hooft 线算子 $T({}^{L}w)$ 在二维有效理论中对 A-膜的作用。经典位形上 $T({}^{L}w)$ 在越过 $y$ 平面时制造指定的奇异：局部 $E_{-}=\oplus_i\mathcal{L}_i$ 变为 $E_{+}=\oplus_i\mathcal{L}_i\otimes\mathcal{O}(z)^{m_i}$，其中 ${}^{L}w=(m_1,\dots,m_N)$ 是 ${}^{L}G$ 的支配权——这正是 Hecke 修正的定义，故 't Hooft 算子等于 Hecke 算子（KW06 §9.1 p.132）。Wilson 线给出谱侧：本征值 $\mathrm{V}_{E}$ 是 ${}^{L}G$-局部系统的表示丛而非数，这与「特征膜」的语言完全一致（见 \ref{core:eigenbrane}）。$S$-对偶交换 't Hooft 与 Wilson，于是把 Hecke 算子与它的谱对象互换。}

\begin{example}[$G=GL_1$：$\mathrm{Pic}$ 上的平移]
取 $G=GL_1$：$\mathrm{Bun}_{GL_1}=\mathrm{Pic}(C)$，对应两侧都是 $\mathrm{Pic}$，$p_z=\mathrm{id}$、$q_z=h_z$，$h_z(L)=L\otimes\mathcal{O}(z)$，故 $H_{1,z}$ 是沿平移的推前（按 Frenkel 2009 §2 的拉回约定写成 $h_z^{*}\mathcal{F}$，两者差一个典范线丛的扭化，不改变特征性）。$\mathrm{Pic}$ 是交换群，故各点的 $H_{1,z}$ 两两对易，可同时对角化——这是特征层在交换情形存在的原因。取 $g=1$：$\mathrm{Pic}^{0}(C)\cong C$ 是复一维环面，$\mathrm{Loc}_{GL_1}(C)=\operatorname{Hom}(\pi_1(C),\mathbb{C}^{\times})\cong(\mathbb{C}^{\times})^{2}$ 是复二维仿射簇。取平凡局部系统，则 $\mathrm{Pic}$ 上带零微分的平凡 $D$-模 $\mathcal{O}_{\mathrm{Pic}}$ 满足 $h_z^{*}\mathcal{O}_{\mathrm{Pic}}\cong\mathcal{O}_{\mathrm{Pic}}$，即特征方程以 $V_{E}|_{z}=\mathbb{C}$ 成立。
\end{example}

\statusline{定义与猜想的出处：Frenkel 2009 §2（Hecke 函子、Hecke 特征层、存在性与唯一性）；KW06 §9.1 p.132（'t Hooft 算子实现 Hecke 修正）、§9.3 p.150（Hecke 对应）、§6.4 p.95（Hecke 特征层的物理定义）。已证：$GL_1$ 由 Deligne 给出（经 Laumon 整理），$GL_2$ 由 Drinfeld，$GL_n$ 由 Frenkel--Gaitsgory--Vilonen（三者见 Frenkel 2009 §2；FGV 的年份出处待核）；任意单群在 oper 情形由 Beilinson--Drinfeld 用二维共形场论构造。$E$ 不可约时特征层的不可约性与唯一性在 Frenkel 2009 §2 中是猜想（$GL_n$ 的已知部分见原文）。函数域的算术原型由 Drinfeld（$GL_2$）与 L.~Lafforgue（一般 $GL_n$）证明（Frenkel 2009 §1），几何对应是它的曲线版本。物理侧的等同（'t Hooft 算子 $=$ Hecke 算子）是 argued：KW06 §9.1 只用经典位形论证，未证明两边的范畴作用严格相同；Witten 2008 §3.1 同样只作物理论证。另一点值得强调：$D$-模必须定义在栈 $\mathrm{Bun}_G$ 上而非半稳定丛的空间上，否则 Hecke 修正会把丛变出定义域、看不见 Hecke 的作用（KW06 §9 与 Witten 2008 §3.1）。}

\dictionaryrow{'t Hooft 线算子与 GNO 分类：\ref{core:line-operators}；特征膜与 Hecke 特征层：\ref{core:eigenbrane}；Hecke 修正空间 $=$ 单极子模空间：\ref{core:monopole-hecke}；$D$-模与扭化：\ref{core:dmod}；整体对应：\ref{core:gl}。}
\end{centry}

\begin{centry}{单极子模空间＝Hecke 修正空间}{Monopole Moduli as Hecke Moduli}{monopole-hecke}
\oneline{带 't Hooft 奇点的三维 Bogomolny 解的模空间，就是曲线上一点处的 Hecke 修正空间；这条等同把「单极子冒泡」与「丛的修正」钉在同一块石头上，也是 't Hooft 算子等于 Hecke 算子的几何根源。}

\begin{definition}[Nahm 数据与奇异 Bogomolny 解]
设 $I=[s_{-},s_{+}]\subset\mathbb{R}$、$k\ge1$。\emph{秩 $k$ 的 Nahm 数据}是三族 $k\times k$ 反 Hermite 矩阵 $T_1,T_2,T_3$（$s\in I$），满足 Nahm 方程
\[
\frac{\mathrm{d}T_i}{\mathrm{d}s}=\frac12\sum_{j,k}\epsilon_{ijk}[T_j,T_k],
\]
并在两端有单极点：$T_i\sim\rho(\sigma_i)/(2(s-s_{\pm}))$，留数构成 $\mathfrak{su}(2)$ 的 $k$ 维不可约表示（此处只用于 $G=SU(2)$）；模去 $U(k)$ 的规范作用 $T_i\mapsto R^{-1}T_iR$。

设 $C$ 为紧黎曼面、$W=I\times C$、$\phi_0\in\Gamma(W,\mathrm{ad}E)$。\emph{奇异 Bogomolny 解}是 $W$ 上 $G$-丛与 $\phi_0$ 满足
\[
F=\star_3D\phi_0 ,
\]
并在某点 $p_0\in C$ 的切片上具有型 ${}^{L}w$ 的 't Hooft 奇点：局部 $E_{-}\cong\oplus_i\mathcal{L}_i$ 越过 $p_0$ 后变为 $E_{+}\cong\oplus_i\mathcal{L}_i\otimes\mathcal{O}(p_0)^{m_i}$。两端取椭圆边界条件：一端固定 $A$（Dirichlet）而 $\phi_0$ 自由，另一端取 $\phi_0=0$ 与协变 Neumann。
\end{definition}

\begin{theorem}[ADHMN 对应]
设 $G=SU(2)$、$k\ge1$。则 $\mathbb{R}^{3}$ 上荷 $k$ 的 BPS 单极子的模空间 $\mathcal{M}_k$ 与秩 $k$ 的 Nahm 数据模空间 $\mathcal{N}_k$ 之间有自然同胚，且该同胚是超凯勒等距，$\dim_{\mathbb{R}}\mathcal{M}_k=4k$。
\end{theorem}
\begin{proof}
ADHMN 构造：给定 Nahm 数据，解一维 Dirac 方程 $D^{\dagger}\psi=0$（$\psi$ 取值于 $\mathbb{C}^{2}\otimes\mathbb{C}^{k}$，$D$ 含 $T_i$），取其 $L^{2}$ 零模 $\psi(s,x)$，置
\[
\phi_0(x)=\int_I s\,\psi^{\dagger}\psi\,\mathrm{d}s,\qquad A_j(x)=\int_I\psi^{\dagger}\partial_j\psi\,\mathrm{d}s ;
\]
对 $x$ 求导并用 Nahm 方程即得 $F=\star_3D\phi_0$。反向由解的 Green 函数在无穷远的展开读出 Nahm 数据。Hitchin 1983 证明该构造穷尽荷 $k$ 的单极子，Donaldson 1984 证明两个模空间的自然度规一致（超凯勒等距）。维数：$k=1$ 时为位置 $3$ 维加整体相位 $1$ 维，共 $4=4k$，一般 $k$ 由指标定理给出。
\end{proof}

\begin{proposition}[单极子模空间即 Hecke 修正空间]
设 $C$ 为紧黎曼面，$p_0\in C$，${}^{L}w$ 为 ${}^{L}G$ 的支配余权，$\mathcal{Y}({}^{L}w;p_0)$ 为 $p_0$ 处型 ${}^{L}w$ 的 Hecke 修正空间（$G$-丛加 $p_0$ 外的同构），$\bar{\mathcal{Y}}$ 为其紧化。则带 $p_0$ 处型 ${}^{L}w$ 的 't Hooft 奇点、两端如上取椭圆边界条件的奇异 Bogomolny 解模空间同构于 $\bar{\mathcal{Y}}({}^{L}w;p_0)$，而 $\bar{\mathcal{Y}}({}^{L}w;p_0)$ 是仿射 Grassmann 流形 $\mathrm{Gr}_{{}^{L}G}$ 中 Schubert 簇 $\bar{Y}({}^{L}w)$ 的闭包。
\end{proposition}
\begin{proof}
取 $y$ 为 $I$ 的坐标。解在 $y=$ 常数的切片上给出 $C$ 上 $G$-丛与 Higgs 场 $(\mathcal{E},\phi)$；越过 $p_0$ 时丛发生一次修正，即上述 $E_{-}\to E_{+}$。反向：给定一次修正，扩展 Bogomolny 方程
\[
\partial_yA_{\bar z}=-\mathrm{i}D_{\bar z}\phi_0,\qquad \partial_y\phi_z=-\mathrm{i}[\phi_z,\phi_0]
\]
沿 $y$ 唯一确定演化（$\phi$ 全纯，故 $p_0$ 外至多一个延拓），于是两侧定义式地重合。真正的定理只是：解空间在奇点外光滑，其维数等于相应椭圆复形的指标；$\phi(p_0)$ 正则半单时只有有限多个允许的修正，由 $\phi$ 在 Schubert 簇上确定的向量场的零点刻画。
\end{proof}

\physside{数学侧「$p_0$ 处一次 Hecke 修正」对应物理侧「一个量子化磁荷单位缩成一点并被 't Hooft 算子吞掉」：$U(1)$-不变的瞬子沿 $\mathbb{C}^{2}/U(1)=\mathbb{R}^{3}$ 的奇点冒泡，权变化 $(m+1,m-1)\to(m,m)$ 就是修正型 ${}^{L}w$ 的退化（KW06 §10.2 p.155）。故 't Hooft 线作用于膜的方式就是对其支撑作 Hecke 修正；而支撑在 Hitchin 纤维上的 type $F$ 膜越过算子后仍支撑在同一纤维，因为 $\phi$ 的特征多项式不随 $y$ 变。这正是 't Hooft 线算子的二维影子等于 Hecke 算子的几何机制（\ref{core:hecke}）。}

\begin{example}[$G=U(2)$：$(1,0)$ 与 $(2,0)$]
取 $G=U(2)$、${}^{L}G=U(2)$。型 ${}^{L}w=(1,0)$：$p_0$ 处的截面只差一个复缩放，修正空间 $\mathcal{Y}\cong\mathbb{CP}^{1}$，紧而光滑（复一维）。型 ${}^{L}w=(2,0)$：未紧化的修正空间是挖去一点的 $\mathbb{P}(1,1,2)$，紧化补入型 $(1,1)$ 的退化修正，得 $\bar{\mathcal{Y}}\cong\mathbb{P}(1,1,2)$，奇点局部为 $\mathbb{C}^{2}/\mathbb{Z}_2$（复二维）。维数对账：指标给出复维数 $a$，$(1,0)$ 对应 $a=1$、$(2,0)$ 对应 $a=2$；对应单极子的荷 $k=a/2$，而荷 $k$ 的单极子模空间实维数 $4k=2a$，恰为修正空间复维数的两倍。
\end{example}

\statusline{ADHMN：Atiyah--Hitchin--Drinfeld--Manin 1978（构造）、Nahm 1979（方程）、Hitchin 1983（$SU(2)$ 的满射性）、Donaldson 1984（模空间等距）。Hecke 修正空间与 Schubert 簇：Beilinson--Drinfeld（具体卷次待核），KW06 §9.2 p.145、§9.3 p.150。命题「单极子模空间 $\cong$ Hecke 修正空间」是 KW06 §10.5 pp.182--183 的构造性等同加物理推导，原文动词为 argued；缺的一步是把「带 't Hooft 奇点的三维 Bogomolny 解」的模空间本身严格定义成数学对象，并证明它就是 Schubert 簇的紧化——原文只把 Schubert 簇的识别当作已引的定理。$\bar{\mathcal{Y}}$ 整体不超凯勒（它来自带边界条件的三维理论），只有奇点附近局部超凯勒。$k=a/2$ 的换算与 $4k$ 的对账是本词典的整理，原文未对账；拿 $(1,0)$ 的 $\mathbb{CP}^{1}$（实二维）直接对 $4k$ 会差因子二，那只是权与荷的换算不同。}

\dictionaryrow{Hecke 修正与算子：\ref{core:hecke}；'t Hooft 线与 GNO 分类：\ref{core:line-operators}；Hitchin 纤维化与 type $F$ 膜：\ref{core:hitchin-fibration}；$\mathcal{M}_H$ 的超凯勒结构：\ref{core:hitchin}。}
\end{centry}

\begin{centry}{D-模与平坦联络（量子 Langlands）}{D-Modules and Flat Connections}{dmod}
\oneline{$D$-模是「带微分算子作用」的层，局部就是平坦联络。几何 Langlands 的一侧正是 $\mathrm{Bun}_G$ 上的（扭化）$D$-模，另一侧是平坦 ${}^{L}G$-丛模空间上的准相干层——这一条把「微分方程」与「谱数据」两种语言接起来。}

\begin{definition}[$D$-模与扭化 $D$-模]
设 $Z$ 为 $\mathbb{C}$ 上光滑代数簇，$\mathcal{D}_Z$ 为由 $\mathcal{O}_Z$ 与向量场生成的层，满足 Leibniz 规则 $[\partial,f]=\partial(f)$。\emph{左 $\mathcal{D}_Z$-模}是带向量场相容作用的 $\mathcal{O}_Z$-模；若该模局部自由，这等价于它上面的可积联络。具体例子：平坦向量丛的截面层是 $D$-模。取线丛 $L$，令 $\mathcal{D}_{L}=\mathcal{D}_Z\otimes_{\mathcal{O}_Z}L$；$\mathcal{D}_{L}$-模称\emph{$L$-扭化 $D$-模}，最基本的例子是 $L$ 的截面层本身（$\mathcal{D}_{L}$ 自然作用在它上面），且 $\mathcal{D}_{L}^{\mathrm{op}}\cong\mathcal{D}_{L^{-1}\otimes K_Z}$（$K_Z$ 为典范丛）。

$\mathrm{Bun}_G$ 是栈而非簇，其上的 $D$-模按 Beilinson--Drinfeld 在导出范畴中定义。几何 Langlands 中出现的扭化有两处：曲线上取 $K_C^{1/2}$ 可平凡化两个等价之间的 $\mathbb{Z}_2$-gerbe（Frenkel 2009 §3 脚注 7）；$\mathrm{Bun}_G$（或其稳定部分的模空间 $M(G,C)$）上取 $K^{1/2}$，使翻译器的弦代数为 $\mathcal{D}_{K^{1/2}}$（KW06 §11.2 p.196，见 \ref{core:bcc}）。
\end{definition}

\begin{theorem}[$D$-模 $=$ 可积联络]
设 $Z$ 为 $\mathbb{C}$ 上光滑簇，$\mathcal{M}$ 为秩 $n$ 的局部自由 $\mathcal{O}_Z$-模。则 $\mathcal{M}$ 上的左 $\mathcal{D}_Z$-模结构与 $\mathcal{M}$ 上的可积联络 $\nabla:\mathcal{M}\to\mathcal{M}\otimes\Omega^{1}_Z$（$\nabla(fs)=f\nabla s+s\,\mathrm{d}f$、$\nabla^{2}=0$）一一对应。特别地，秩一 $D$-模局部由闭形式 $\eta\in\Omega^{1}_Z$ 的 $\nabla=\mathrm{d}+\eta$ 刻画，同构类由「单值」分类。
\end{theorem}
\begin{proof}
$\mathcal{D}_Z$ 由 $\mathcal{O}_Z$ 与 $\mathrm{Der}(Z)$ 生成，关系就是 $[\partial,f]=\partial(f)$。故给 $\mathcal{D}_Z$-模结构等价于给 $\mathcal{O}_Z$-模结构再加向量场的作用 $\partial\mapsto\nabla_{\partial}$；后者与函数乘法相容，当且仅当 $\nabla_{\partial}(fs)=f\nabla_{\partial}s+\partial(f)s$，即 $\nabla$ 是联络。两个向量场作用的交换子必须对应 $[\partial,\partial']$ 的作用，这恰是曲率为零：$[\nabla_{\partial},\nabla_{\partial'}]=\nabla_{[\partial,\partial']}$，即 $\nabla^{2}=0$。反向沿同一对应逐项检查即得。
\end{proof}

\physside{数学侧「$M(G,C)$ 上的 $D$-模」对应物理侧「$M(G,C)$ 的余切丛 $T^{*}M(G,C)$ 的非交换形变上的 B-膜」：Kapustin 的论证是，全纯辛流形上的 A-膜范畴与其非交换形变上的 B-膜范畴等价，而 $T^{*}M(G,C)$ 的非交换形变正是把 A-模型（只看辛形式）与 B-模型（只看复结构）接起来的对象（KW06 §11.1 pp.185--189）。又因 $\mathcal{M}_H(G,C)$ 以 $T^{*}M(G,C)$ 的稳定部分为 Zariski 开集，$\mathcal{M}_H$ 上的 A-膜可经典范余各向同性膜翻译成这些 $D$-模。于是「$D$-模 $=$ 平坦联络」这条数学事实的物理读法是：平坦联络就是余切丛上 B-膜的经典极限。}

\begin{example}[$\mathbb{C}^{\times}$ 与 $\mathrm{Bun}_{GL_1}$]
取 $Z=\mathbb{C}^{\times}$：对 $\lambda\in\mathbb{C}$，$\nabla_{\lambda}=\mathrm{d}+\lambda\,\mathrm{d}z/z$ 给出秩一 $D$-模；$\lambda-\lambda'\in\mathbb{Z}$ 时两者同构（用 $z^{k}$ 作规范变换），否则不同构，其单值为 $\exp(2\pi\mathrm{i}\lambda)$——故秩一 $D$-模由单值分类，这是「$D$-模 $=$ 局部系统」的最小验证。取 $G=GL_1$、$C$ 亏格 $g$：$\mathrm{Bun}_{GL_1}=\mathrm{Pic}(C)$ 复维数 $g$，而 $\mathrm{Loc}_{GL_1}(C)\cong(\mathbb{C}^{\times})^{2g}$ 复维数 $2g$；两侧维数不等并不矛盾——对应发生在导出范畴的层次，且 $\mathrm{Loc}_{GL_1}$ 与 $T^{*}\mathrm{Pic}^{0}$ 实维数相同，因为 $\mathcal{M}_H(GL_1)$ 在这两个复结构下分别是它们。
\end{example}

\statusline{$D$-模与扭化 $D$-模、$\mathcal{D}_{L}^{\mathrm{op}}\cong\mathcal{D}_{L^{-1}\otimes K}$：KW06 §11.2 pp.189--200（p.196 处），定义与平坦联络的联系见 Frenkel 2009 §2。范畴等价 $D^{b}(\mathcal{D}\text{-mod}(\mathrm{Bun}_G))\simeq D^{b}(\mathrm{QCoh}(\mathrm{Loc}_{{}^{L}G}))$：Frenkel 2009 §3 式 (3.1)，原文列为猜想；$GL_1$ 已由 Laumon、Rothstein 证明（Fourier--Mukai），$GL_n$ 的 Hecke 特征层由 Deligne、Drinfeld、Frenkel--Gaitsgory--Vilonen 构造，完整等价仍是猜想（据 Frenkel 2009；近年 Gaitsgory 与合作者的系列论文宣称证明一般约化群情形，卷次与年份待核）。Kapustin 的「$D$-模 $=$ 非交换形变上的 B-膜」是物理推导（KW06 §11.1 p.185 原文用 argued；Witten 2008 §3.1 同）。严格可对照的结果是 Nadler--Zaslow 2009 关于 $T^{*}M$ 的 Fukaya 范畴与 $M$ 上可构造层的等价（细节待核），那不直接给出本条目的等价。}

\dictionaryrow{平坦联络与局部系统、对偶群：\ref{core:lg}；翻译器与 $\mathcal{D}_{K^{1/2}}$：\ref{core:bcc}；Hecke 特征层：\ref{core:hecke}；几何 Langlands 对应：\ref{core:gl}；$\mathcal{M}_H$ 的复结构：\ref{core:hitchin}。}
\end{centry}

\begin{centry}{几何 Langlands 对应}{The Geometric Langlands Correspondence}{gl}
\oneline{$\mathrm{Bun}_G$ 上的 $D$-模与 $\mathrm{Loc}_{{}^{L}G}$ 上的对象之间的等价：Weil 三栏石碑里「黎曼面」那一栏的现代形态；物理上由 $S$-对偶、紧化与镜像对称三步推出。}

\begin{definition}[非分歧几何 Langlands 对应]
设 $C$ 为 $\mathbb{C}$ 上光滑射影连通曲线，$G$ 为复约化群，${}^{L}G$ 为其对偶群。记 $\mathrm{Bun}_G$ 为 $C$ 上 $G$-丛的模栈，$\mathrm{Loc}_{{}^{L}G}(C)$ 为 $C$ 上平坦 ${}^{L}G$-丛（即局部系统 $\pi_{1}(C)\to{}^{L}G$）的模栈。\emph{非分歧几何 Langlands 对应}断言：对每个 $E\in\mathrm{Loc}_{{}^{L}G}(C)$ 存在非零 $D$-模 $\mathcal{F}_{E}$ 满足 Hecke 特征方程
\[
H_{V,z}(\mathcal{F}_{E})\cong V_{E}|_{z}\otimes\mathcal{F}_{E}\qquad(z\in C,\ V\in\operatorname{Rep}{}^{L}G);
\]
\emph{范畴化形式}断言导出范畴的等价
\[
D^{b}\bigl(\mathcal{D}\text{-mod}(\mathrm{Bun}_G)\bigr)\;\simeq\;D^{b}\bigl(\mathrm{QCoh}(\mathrm{Loc}_{{}^{L}G}(C))\bigr),\qquad \mathcal{O}_{E}\mapsto\mathcal{F}_{E}.
\]
在 $C$ 上标若干点并指定局部数据（正则奇点、抛物结构等）时称\emph{分歧}情形；本词典与 KW06 一样只处理非分歧情形。
\end{definition}

\begin{theorem}[$GL_1$：几何类域论]
设 $C$ 为 $\mathbb{C}$ 上光滑射影连通曲线，$G=GL_1$。则上述范畴等价成立：
\[
D^{b}\bigl(\mathcal{D}\text{-mod}(\mathrm{Pic}(C))\bigr)\;\simeq\;D^{b}\bigl(\mathrm{QCoh}(\mathrm{Loc}_{GL_1}(C))\bigr).
\]
\end{theorem}
\begin{proof}
构造是 Fourier--Mukai 型变换：取 $\mathrm{Pic}(C)\times\mathrm{Loc}_{GL_1}(C)$ 上的 Poincar\'e 对象为核（沿 $\mathrm{Pic}$ 方向是平坦的 $D$-模、沿 $\mathrm{Loc}$ 方向是准相干的），置 $\mathcal{F}\mapsto (p_{\mathrm{Loc}})_{*}\bigl(p_{\mathrm{Pic}}^{*}\mathcal{F}\otimes\mathcal{P}\bigr)$。核的平移相容性给出 $h_z^{*}\mathcal{P}\cong\mathcal{P}\otimes\mathbb{E}_{z}$（$\mathbb{E}$ 为 $\mathrm{Loc}$ 上的泛局部系统），于是 $\mathrm{Loc}$ 上的摩天大楼层 $\mathcal{O}_{E}$ 被送到满足 $H_{1,z}\mathcal{F}\cong E|_{z}\otimes\mathcal{F}$ 的 $D$-模，正是 Hecke 特征层；逆变换由核的对偶给出。Laumon、Rothstein 证明该变换是导出范畴的等价（Frenkel 2009 §3；核的具体构造与整性条件此处从略）。
\end{proof}

\physside{数学侧「$\mathrm{Bun}_G$ 上的 $D$-模 $=$ 平坦 ${}^{L}G$-丛的谱数据」对应物理侧「四维 $\mathcal{N}=4$ 理论在 $C$ 上紧化后，两种 $S$-对偶扭化给出的镜像对」。推导链（KW06 §4--§5、§11）：把理论放在 $\Sigma\times C$ 上，在 $C$ 上紧化得靶为 $\mathcal{M}_H(G,C)$ 的二维拓扑 sigma 模型；$(G_c,t=1,\theta=0)$ 给出 $\mathcal{M}_H(G)$ 上辛形式 $\omega_K$ 的 A-模型，$({}^{L}G_c,{}^{L}t=\mathrm{i},{}^{L}\theta=0)$ 给出 $Y({}^{L}G)=\mathcal{M}_H({}^{L}G)$ 复结构 $J$ 的 B-模型，$S$-对偶断言两者相等价，即镜像对称（Hitchin 纤维化的底相同、纤维互为对偶环面）。再用 (i) 同调镜像对称：$Y({}^{L}G)$ 上 B-膜 $\simeq\mathcal{M}_H(G)$ 上 A-膜；(ii) 典范余各向同性膜：$\mathcal{M}_H(G)$ 上 A-膜 $\simeq\mathrm{Bun}_G$ 上扭化 $D$-模；复合两步即得对应。Hecke 算子对应 ${}^{L}G$ 局部系统的表示丛 $V_{E}$：$S$-对偶把 't Hooft 线（即 Hecke 算子）变成 Wilson 线，而 Wilson 线的本征值就是 $V_{E}$。}

\begin{example}[$G=GL_1$ 与 Hitchin 模空间]
取 $G=GL_1$、$C$ 亏格 $g$：一侧 $\mathrm{Pic}(C)$ 复维数 $g$（含 $\mathbb{Z}$ 个分支），另一侧 $\mathrm{Loc}_{GL_1}(C)\cong(\mathbb{C}^{\times})^{2g}$ 复维数 $2g$。$g=1$ 时 $\mathrm{Pic}^{0}(C)\cong C$ 为一维环面、$\mathrm{Loc}\cong(\mathbb{C}^{\times})^{2}$，对应就是 $C$ 上的 Fourier--Mukai 变换（见 \ref{core:dmod} 的例）。关键在于两侧为何能对上：$\mathcal{M}_H(GL_1,C)$ 在复结构 $I$ 下等于 $T^{*}\mathrm{Pic}^{0}(C)$（复维数 $2g$），在复结构 $J$ 下等于 $\mathrm{Loc}_{GL_1}(C)$（复维数 $2g$）——同一个超凯勒流形的两个复结构，故「$S$-对偶归结为 $T$-对偶」在这里可以逐维核对。
\end{example}

\statusline{Weil 的三栏石碑（1940 年致 Simone Weil 的信：数论、$\mathbb{F}_q$ 上的曲线、黎曼面三栏；原文有「我的工作就是解读一篇三语碑文」之语）：Frenkel 2009 §1，引 Weil, Notices AMS 52 (2005) 334--341。陈述、猜想与范畴化：Frenkel 2009 §2--§3。已证：$GL_1$（Laumon--Rothstein 的 Fourier--Mukai）、$GL_2$（Drinfeld）、$GL_n$ 的特征层存在性（Frenkel--Gaitsgory--Vilonen，年份待核）、oper 情形（Beilinson--Drinfeld）；一般群的完整范畴等价仍是猜想（近年 Gaitsgory 与合作者的系列论文宣称证明，卷次待核）。物理推导（$S$-对偶 $+$ 紧化 $+$ 镜像 $+$ 翻译器）：KW06 §4--§5、§11 与 Witten 2008 §3.1、§5.3，原文动词为 argued；缺的是把四维与二维的膜范畴严格定义并证明两步等价——同调镜像对称的 B-膜一侧与典范余各向同性膜的翻译器都只有物理论证，唯一严格的替代是 Nadler--Zaslow 对 $T^{*}M$ 的构造（细节待核）。分歧情形见 KW06 §4（页码待核）与 Frenkel--Gaitsgory、Gukov--Witten，本词典不展开。}

\dictionaryrow{对偶群与 Langlands 纲领：\ref{core:lg}；$D$-模：\ref{core:dmod}；Hecke 算子与特征层：\ref{core:hecke}；翻译器 $B_{cc}$：\ref{core:bcc}；$S$-对偶与镜像对称：\ref{core:em-duality}、\ref{core:ab-model}。}
\end{centry}

\begin{centry}{翻译器 $B_{cc}$}{The Canonical Coisotropic Brane}{bcc}
\oneline{在 $\mathcal{M}_H(G,C)$ 上空间充满的典范余各向同性 A-膜：它的自弦代数恰是微分算子层，于是「A-膜 $\mapsto$ $\mathrm{Bun}_G$ 上扭化 $D$-模」这一翻译有了载体——这是几何 Langlands 物理推导的最后一步。}

\begin{definition}[典范余各向同性 A-膜]
辛流形 $(M,\omega)$ 上空间充满、Chan-Paton 为线丛 $\mathcal{N}$（联络 $\alpha$、曲率 $F=\mathrm{d}\alpha$）的 A-膜条件为 $(\omega^{-1}F)^{2}=-1$（余各向同性膜，见 \ref{core:abrane}）。取 $M=\mathcal{M}_H(G,C)$、复结构 $K$ 对应的辛形式 $\omega_{K}$，令
\[
\omega=(\operatorname{Im}\tau)\,\omega_{K},\qquad F=(\operatorname{Im}\tau)\,\omega_{J},
\]
其中 $\omega_{J}$ 是同一超凯勒结构中的另一个辛形式。\emph{典范余各向同性膜} $\mathcal{B}_{cc}$ 就是 $M$ 上取此 $F$ 的秩一 A-膜，支撑为整个 $M$（$M$ 以 $T^{*}\mathrm{Bun}_G$ 的稳定部分为 Zariski 开集，故也可说它支撑在 $\mathrm{Bun}_G$ 的余切丛上）。$F$ 恰当：$F=\mathrm{d}\alpha$，且 $\alpha$ 可由规范场论显式给出
\[
\alpha=\frac{\operatorname{Im}\tau}{2\pi}\int_{C}|\mathrm{d}^{2}z|\;\operatorname{Tr}\bigl(\phi_{z}\,\delta A_{\bar z}+\phi_{\bar z}\,\delta A_{z}\bigr),
\]
故 $\mathcal{N}$ 的整体存在性等价于 $\omega_{J}$ 的周期是 $2\pi$ 的整数倍（KW06 §11.2 式 (11.3)）；同族的 $(\operatorname{Im}\tau)(a\omega_{J}+b\omega_{I})$（$a^{2}+b^{2}=1$）在 $b\neq0$ 时不是任何联络的曲率。$G$ 单连通时 $\mathcal{B}_{cc}$ 唯一，否则唯一到扭化平坦线丛。
\end{definition}

\begin{theorem}[$\mathcal{B}_{cc}$ 满足 A-膜条件]
设 $\mathcal{M}_H(G,C)$ 在相应点光滑，$\omega=(\operatorname{Im}\tau)\omega_{K}$、$F=(\operatorname{Im}\tau)\omega_{J}$。则：
\begin{enumerate}[label=(\roman*)]
\item $(\omega^{-1}F)^{2}=-1$，且 $\hat{N}=\omega^{-1}F$ 是 $\mathcal{M}_H$ 的复结构 $I$；
\item $\omega-\mathrm{i}F$ 关于 $\hat{N}=I$ 是 $(2,0)$ 型全纯辛形式，故 $\hat{N}$ 可积；
\item $F$ 闭；$\mathcal{B}_{cc}$ 整体存在时其 de Rham 类恰为 $2\pi$ 的整数倍。
\end{enumerate}
\end{theorem}
\begin{proof}
(i) $\omega^{-1}F=(\operatorname{Im}\tau)^{-1}(\operatorname{Im}\tau)\,\omega_{K}^{-1}\omega_{J}=\omega_{K}^{-1}\omega_{J}$。由超凯勒的 $K=IJ$ 得 $\omega_{J}(u,v)=g(Ju,v)=g(IJIu,v)$ 与 $\omega_{K}(Iu,v)=g(KIu,v)=g(Ju,v)$，故 $\omega_{K}^{-1}\omega_{J}=I$，而 $I^{2}=-1$。
(ii) 在复结构 $I$ 下 $\Omega_{I}=\omega_{J}+\mathrm{i}\omega_{K}$ 是全纯辛 $(2,0)$ 形式，故 $\omega-\mathrm{i}F=(\operatorname{Im}\tau)(\omega_{K}-\mathrm{i}\omega_{J})=-\mathrm{i}(\operatorname{Im}\tau)\Omega_{I}$ 是 $(2,0)$ 型且全纯。
(iii) $\mathrm{d}F=(\operatorname{Im}\tau)\mathrm{d}\omega_{J}=0$；$\mathrm{d}(\omega-\mathrm{i}F)=0$ 给出 $\hat{N}$ 的可积性（KW06 §11.2 引 ref80）。整性来自 $\mathcal{N}$ 的整体存在性。
\end{proof}

\begin{proposition}[翻译器]
设 $M=\mathcal{M}_H(G,C)$ 光滑，$\mathcal{B}_{cc}$ 如上。则 $A\mapsto\operatorname{Hom}(\mathcal{B}_{cc},A)$ 给出 $M$ 上复结构 $K$ 的 A-膜范畴到 $\mathrm{Bun}_G$ 上扭化 $D$-模导出范畴的等价。扭化来自时间反演：$\mathcal{B}_{cc}$ 在时间反演下不变，故其开弦代数须同构于自己的相反代数；由 $\mathcal{D}_{L}^{\mathrm{op}}\cong\mathcal{D}_{L^{-1}\otimes K_{M}}$ 得 $L^{2}\cong K_{M}$，即 $L=K_{M}^{1/2}$。两个平方根相差 $K_{M}$ 的 $2$-挠线丛，即相差一个 $\mathbb{Z}_{2}$-gerbe 的平凡化，无典范选择。
\end{proposition}
\begin{proof}
要点四条。(i) 对任意两膜 $A,A'$，$(A,A')$ 弦构成 $(A,A)$ 弦代数的模（开--闭拓扑场论公理），故取 $A=\mathcal{B}_{cc}$ 得到 $\mathcal{D}_{K^{1/2}}$ 的一个模。(ii) 取 $A$ 为支撑在 $\phi=0$ 的 $M(G,C)$ 拷贝上的 $(B,A,A)$ 膜：法向费米子两端反号，无零模，量子化给出旋量丛 $K_{M}^{1/2}$ 的截面，超荷为 $\bar\partial$，故 $\operatorname{Hom}(\mathcal{B}_{cc},A)$ 是 $K_{M}^{1/2}$ 的全纯截面层——这解释了自弦代数为何恰是 $\mathcal{D}_{K^{1/2}}$。(iii) 圆盘上的瞬子被边界解析性逼成常数，故无瞬子修正，微扰部分按 $\Omega_{I}$ 的 Poisson 括号非交换化。(iv) 要落到 $\mathrm{Bun}_G$ 而非 $M(G,C)$，必须把层搬到栈上：Hecke 修正会把稳定丛变出稳定部分。
\end{proof}

\physside{数学侧「$M$ 上复结构 $K$ 的 A-膜」对应物理侧「四维理论中保持相应超荷的边界条件」，而 $\mathcal{B}_{cc}$ 是其中空间充满的那一条；$\operatorname{Hom}(\mathcal{B}_{cc},A)$ 是「一端贴 $\mathcal{B}_{cc}$、一端贴 $A$ 的开弦态空间」，其上的微分算子作用来自边界作用量 $\mathcal{S}=-\mathrm{i}\operatorname{Im}\tau\int_{\partial\Sigma}\Phi^{*}\varpi$ 的量子化。$\mathcal{B}_{cc}$ 是 $(A,B,A)$ 型：在 $J$ 中它是 B-膜（支撑为全空间、$F$ 与 $\omega_{J}$ 成比例，故 Chan-Paton 曲率为 $J$ 型的 $(1,1)$ 形式）；在 $I$ 与 $K$ 中则是 A-型的余各向同性膜，依据是恒等式 $(\omega_{I}^{-1}\omega_{J})^{2}=-1$（KW06 §11.2 式 (11.5)）。§12 把这套构造装回四维：$(A,B,A)$ 型膜的边界条件由 $\Gamma$ 矩阵 $T'=-\Gamma_{15}\Gamma_{2-7}\Gamma_{3+6}$ 给出，其对偶 $\hat{T}$ 的混合 Dirichlet--Neumann 条件正好读出 $f=\omega_{J}$，即 $\mathcal{B}_{cc}$（KW06 §12.3 pp.215--218）。}

\begin{example}[$G=SU(2)$、$g=2$]
取 $C$ 亏格 $2$、$G=SU(2)$：$\dim_{\mathbb{C}}\mathrm{Bun}_G=3$，故 $\dim_{\mathbb{C}}\mathcal{M}_H=6$。$\mathcal{B}_{cc}$ 支撑整个 $\mathcal{M}_H$，其 $F=(\operatorname{Im}\tau)\omega_{J}$；由定理 (i) 得 $\hat{N}=I$，与 $\tau$ 无关，故对任意复化耦合都存在。$\mathcal{M}_H$ 超凯勒，全纯辛形式 $\Omega_{I}$ 给出 $K_{M}\cong\mathcal{O}_{M}$ 的处处非零截面，故 $K_{M}^{1/2}$ 是 $2$-挠线丛——这正是 $\mathbb{Z}_{2}$-gerbe 不平凡但不障碍等价的具体形态。取 $G=U(1)$、$g=1$：$\mathcal{M}_H\cong T^{*}\mathrm{Pic}^{0}\cong T^{2}\times\mathbb{C}$ 复二维，$\mathcal{B}_{cc}$ 的自弦代数 $=\mathcal{D}_{K^{1/2}}$，而在 $\mathbb{C}^{2}$ 局部上它就是 Weyl 代数，与「余切丛非交换形变上的 B-膜」一致。
\end{example}

\statusline{定义、条件 $(\omega^{-1}F)^{2}=-1$、$F=(\operatorname{Im}\tau)\omega_{J}$、显式 $\alpha$：KW06 §11.1 p.185、§11.2 pp.189--200（式 (11.1)--(11.3)）。定理 (i)--(iii) 是 KW06 §11.2 的逐项计算，可独立核验（本条目的证明）。命题（$\operatorname{Hom}(\mathcal{B}_{cc},\cdot)$ 是等价）是 argued：KW06 §11.1 p.185 原文用 argued 并两处说 roughly speaking（严格对象是栈 $\mathrm{Bun}_G$ 上的扭化 $D$-模，且须选合适的 A-膜范畴），§11.3 pp.200--206 只对 $GL_1$ 与 $U(1)$ 情形作出可核对的计算。缺的一步是：(a) 把「含余各向同性对象的 A-膜范畴」与开弦代数严格定义，且层化只在 $T^{*}\mathcal{V}$ 型开集上可行；(b) 证明 $\operatorname{Hom}$ 函子是范畴等价；(c) 把 $M(G,C)$ 上的层严格搬到 $\mathrm{Bun}_G$ 栈上。已知严格特例：$GL_1$ 的 Fourier--Mukai（Laumon--Rothstein）、$GL_n$ 的特征层、oper（Beilinson--Drinfeld），以及 Nadler--Zaslow 对 $T^{*}M$ 的构造（细节待核）。§11.4（约 p.206，页码待核）把 $\Psi$ 推广：$F=(\operatorname{Im}\tau)\cos q\,\omega_{J}$、$B=-\operatorname{Re}\tau\,\omega_{I}$，弦代数为 $\mathcal{D}_{K_{M}^{1/2}\otimes\mathfrak{L}^{\Psi}}$；$\Psi=h$（$h$ 为对偶 Coxeter 数）时化为未扭化 $D$-模，再与 $S$-对偶复合即得几何 Langlands 的标准表述。}

\dictionaryrow{余各向同性 A-膜与 A-膜条件：\ref{core:abrane}；$D$-模与扭化：\ref{core:dmod}；几何 Langlands 对应：\ref{core:gl}；$\mathcal{M}_H$ 的超凯勒结构与 $\omega_{I},\omega_{J},\omega_{K}$：\ref{core:hitchin}；四维边界条件与膜类型：\ref{core:ab-model}。}
\end{centry}
